%
% 6.006 problem set 0 solutions template
%
\usepackage{amsthm}

\documentclass[12pt,twoside]{article}

\input{macros-sp20}
\newcommand{\theproblemsetnum}{0}

\title{6.006 Problem Set 0}

\begin{document}

\handout{Problem Set \theproblemsetnum}

\setlength{\parindent}{0pt}
\medskip\hrulefill\medskip

{\bf Name:} Nicolas Iskandar

\medskip\hrulefill

\begin{problems}

\problem

\begin{problemparts}
\problempart \{6, 12\}
\problempart 7
\problempart 3
\end{problemparts}

\problem

\begin{problemparts}
\problempart $ E[x] = E[x_1] + E[x_2] + E[x_3] \\
E[x_i] = \frac{1}{2} * 1 + \frac{1}{2} * 0 = \frac{1}{2} \\
E[x] = \frac{3}{2} $
\problempart $ E[Y] = E[D_1] * E[D_2] \\
E[D_i] = \frac{1}{6}(1+2+3+4+5+6) = \frac{21}{6} = 3.5\\
E[Y] = 3.5  * 3.5 = 12.25 $
\problempart $ E[X+Y] = E[X] + E[Y] = 1.5 + 12.25 = 13.75 $
\end{problemparts}

\problem

\begin{problemparts}
  $ A=\frac{600}{6}=100$ \\ $ B = 60\mod42 = 18 $
\problempart True
\begin{align*}
100 \mod 2 &= 0 \\
 18 \mod 2  &= 0 \\
\implies A &\equiv B \pmod{2}
\end{align*}
\problempart False
\begin{align*}
100 \mod 3 &= 1 \\
 18 \mod 3  &= 0 \\
\implies A &\not\equiv B \pmod{3}
\end{align*}
\problempart False
\begin{align*}
100 \mod 4 &= 0 \\
 18 \mod 4  &= 2 \\
\implies A &\not\equiv B \pmod{4}
\end{align*} 
\end{problemparts}

\problem \\
Base case: for $n = 1$,
\begin{align*}
\sum_{i=1}^{1}i^3 &= [\frac{1*(2)}{2}]^2 \\
1 &= 1 \implies true
\end{align*}
Inductive Step: Suppose the equation is true for $n=k$, prove that it's true for $n=k+1$
\begin{align*}
\sum_{i=1}^{k+1} i^3
  &= (k+1)^3 + \sum_{i=1}^{k} i^3 \\
  &= (k+1)^3 + \left(\frac{k(k+1)}{2}\right)^2 \\
  &= (k+1)^3 + \frac{k^2 (k+1)^2}{4} \\
  &= \frac{4(k+1)^3 + k^2 (k+1)^2}{4} \\
  &= \frac{(k+1)^2 (k^2 + 4k + 4)}{4}
   = \left(\frac{(k+1)(k+2)}{2}\right)^2 \qedbox
\end{align*}

\newpage
\problem \\
Base case: $k = 1$.
$\implies$ The graph has 1 vertex and $|E| = 0$ edges
$\implies$ A graph with one vertex and no edges is trivially acyclic 
$\implies$ Therefore, the base case holds.

Inductive hypothesis: Assume that every connected graph with $k$ vertices and $k-1$ edges is acyclic.

Inductive step: Consider a connected graph $G$ with $k+1$ vertices and $k$ edges.

Let $\deg(v)$ denote the degree of vertex $v$.
The sum of degrees of all vertices is $2|E| = 2k$. \\ 
The average degree is \\
\[\frac{\sum_{v \in V} \deg(v)}{|V|} = \frac{2k}{k+1} < 2.\]
Since the average degree is less than 2, there exists at least one vertex $v$ with degree 1. Let $v$ be connected to a single vertex $u$.
Remove vertex $v$ and its single edge $(v,u)$. Denote the remaining graph as $G_0$.\\
$G_0$ has $k$ vertices and $k-1$ edges.
$G_0$ is connected because removing a leaf cannot disconnect the graph.
By the inductive hypothesis, $G_0$ is acyclic.\\
Any cycle in $G$ would have to include $v$, but $\deg(v) = 1$, so $v$ cannot be part of any cycle.\\
Hence, $G$ is acyclic \qedbox

\problem
\begin{lstlisting}
def count_long_subarray(A):
    '''
    Input:  A     | Python Tuple of positive integers
    Output: count | number of longest increasing subarrays of A
    '''
    count = 0
    i = 0
    j = 0
    maxlength = 0
    n = len(A)
    
    while j < n:
      while j < n - 1 and A[j + 1] > A[j]:
        j += 1
        
      length = j - i + 1
      if length > maxlength:
        maxlength = length
        count = 1
      elif length == maxlength:
        count += 1
      
      j += 1
      i = j
    
    return count
\end{lstlisting}

\end{problems}

\end{document}
